% !TEX root = ../main.tex
% !TEX program = lualatex
\documentclass[../main.tex]{subfiles}
\IfSubfilesClassLoaded{\externaldocument{\subfix{../build/main}}}{}
% =============================================================================
% CHAPTER 12: SOLICITATION
% DESCRIPTION: Solicitation prep steps, practitioner checklists, and Navy overlays.
% =============================================================================
\begin{document}
\IfSubfilesClassLoaded{\chapter{Solicitation Preparation}}{}

% === SOLICITATION PREPARATION (EDO 3.2.2) ===
\section{Solicitation Preparation}

\subsection{Summary}
\begin{itemize}
  \item The two primary contract categories are fixed-price and cost-reimbursement.
  \item Fixed-price contracts transfer cost risk to the contractor.
  \item An \ac{idiq} contract provides for indefinite quantities within stated limits during a fixed period.
  \item Pre-solicitation sequence: determination of need, requirements package, and method of contracting.
  \item \ac{ifb} is used for sealed bidding when conditions for sealed bidding are met.
\end{itemize}
~\srcCite{edo-3-2-2-solicitation-prep-2025,FAR-16-103-Negotiating-Contract-Type,FAR-16-202-1-FFP,FAR-16-301-1-Cost-Reimbursement,FAR-16-501-2-Indefinite-Delivery}

\subsection{Contracting Process Overview}
\begin{itemize}
  \item Pre-solicitation phase (determination of need, requirements package, method of contracting).
  \item Solicitation, evaluation, and award phase (solicitation of sources, response and evaluation, award).
  \item Post-award phase (contract administration, delivery, closeout).
\end{itemize}
~\autocite{edo-3-2-2-solicitation-prep-2025}

\subsection{Pre-Solicitation Phase}
\begin{description}
  \item[\textbf{Determination of capability need.}] For systems with new capabilities, the contracting process starts only after a need is identified, non-materiel solutions are ruled out, and approval is granted to pursue a materiel solution~\autocite{edo-3-2-2-solicitation-prep-2025}.
  \item[\textbf{Requirements package.}] The package establishes the baseline for solicitation development and includes~\autocite{edo-3-2-2-solicitation-prep-2025}:
        \begin{itemize}
          \item Management reviews and approvals.
          \item Certification of funds availability.
          \item Requirements description (specifications, \ac{sow}/\ac{soo}, purchase descriptions, contract data requirements).
          \item Special packaging and marking requirements.
          \item Recommended evaluation factors and criteria (team-developed).
          \item Inspection and acceptance requirements.
          \item Delivery or performance schedule.
          \item Special contract administration requirements.
          \item Special provisions or clauses (e.g., options, \ac{gfe}).
          \item Recommended potential sources and market research results.
          \item Input for other-than-full-and-open competition approvals.
          \item Acquisition planning documentation required by the current acquisition and agency rules; there is no universal \$10M action-level test.
        \end{itemize}
\end{description}
\note{A contract requirements package requires substantial effort across the team~\autocite{edo-3-2-2-solicitation-prep-2025}.}

\note{Acquisition planning applies broadly. The current \ac{rfo} Part~7 calls for written plans for cost-reimbursement and other high-risk contracts. The published DoD Part~7 baseline also uses program-level tests: development of \$10M or more; production or services of \$50M across all fiscal years or \$25M in one fiscal year; and other programs as appropriate, with specified exceptions. Confirm the signed DoD/Navy deviation and approval rules for the actual action~\autocite{FAR-Overhaul-Part7-2026,DoD-RFO-Deviation-Part7-2026}.}

\subsection{Statement of Objectives and Performance Work Statement}
\begin{description}
  \item[\textbf{SOO.}] A Government-prepared, high-level objectives document that is provided in the solicitation in lieu of a Government-written \ac{pws}; offerors use it to propose the \ac{pws}~\autocite{edo-3-2-2-solicitation-prep-2025}.
  \item[\textbf{PWS.}] Describes work in terms of required results rather than how the work is performed or hours provided; may be Government-developed or derived from an \ac{soo}~\autocite{edo-3-2-2-solicitation-prep-2025}.
\end{description}
\note{Example: \ac{soo} = ``maintain attractive, functional groundcover''; \ac{pws} = ``maintain grass between 1 and 2 inches.''~\autocite{edo-3-2-2-solicitation-prep-2025}.}

\subsection{Pre-Solicitation Communication and Market Research}
\begin{itemize}
  \item Before soliciting offers, the requirement is advertised to gauge industry interest through market research, industry days, \acp{rfi}, and draft \acp{rfp}~\autocite{edo-3-2-2-solicitation-prep-2025}.
  \item Robust communication improves the quality of the \ac{rfp} and the source selection plan by clarifying Government needs and industry capabilities~\autocite{edo-3-2-2-solicitation-prep-2025}.
\end{itemize}

\subsection{Market Research Purpose}
\begin{itemize}
  \item Conducted throughout the acquisition life cycle to inform acquisition strategy and planning~\autocite{edo-3-2-2-solicitation-prep-2025}.
  \item Evaluates commercial marketplace capability and identifies potential sources~\autocite{edo-3-2-2-solicitation-prep-2025}.
  \item Supports price analysis and scales in depth based on urgency, dollar value, complexity, and past experience~\autocite{edo-3-2-2-solicitation-prep-2025}.
  \item Conducted in an \ac{ipt} environment~\autocite{edo-3-2-2-solicitation-prep-2025}.
\end{itemize}

\subsection{IPT Roles in Market Research}
\begin{description}
  \item[\textbf{PM.}] Owns overall responsibility, develops the acquisition strategy, and ensures transparency in sources-sought and market research activities~\autocite{edo-3-2-2-solicitation-prep-2025}.
  \item[\textbf{KO.}] Negotiates, awards, administers, and terminates contracts; determines whether the requirement can be set aside for small business; supported by contract specialists, cost analysts, and other functional experts~\autocite{edo-3-2-2-solicitation-prep-2025}.
  \item[\textbf{Technical IPT.}] Evaluates existing commercial products and non-developmental items for applicability~\autocite{edo-3-2-2-solicitation-prep-2025}.
  \item[\textbf{Cost analysts.}] Provide input on appropriate contract pricing information~\autocite{edo-3-2-2-solicitation-prep-2025}.
\end{description}

\subsection{Methods of Contracting}
\begin{description}
  \item[\textbf{Sealed bidding.}] Used when award can be made on price, specifications are clear, discussions are unnecessary, and there is a reasonable expectation of more than one bid. Uses an \ac{ifb} and awards to the lowest priced, responsive, responsible bidder~\autocite{edo-3-2-2-solicitation-prep-2025}.
  \item[\textbf{Negotiation.}] Used when any sealed-bidding condition is lacking (time, price information, requirements definition, or competition). May be sole source or competitive using an \ac{rfp}; awards are based on evaluation criteria and discussions as needed~\autocite{edo-3-2-2-solicitation-prep-2025}.
\end{description}

\subsection{Contract Types Overview}
\begin{itemize}
  \item Contract type is a key pre-solicitation decision tied to risk, requirements clarity, and profit incentives~\autocite{edo-3-2-2-solicitation-prep-2025,FAR-16-103-Negotiating-Contract-Type}.
  \item Two broad categories: fixed-price and cost-reimbursement (cost-plus)~\autocite{edo-3-2-2-solicitation-prep-2025,FAR-16-102-Contract-Type-Policies,FAR-16-202-1-FFP,FAR-16-301-1-Cost-Reimbursement}.
\end{itemize}

\subsection{Fixed-Price Contracts}
\begin{itemize}
  \item Provide for a firm price or adjustable price constrained by a ceiling or target price~\autocite{edo-3-2-2-solicitation-prep-2025,FAR-16-202-1-FFP,FAR-16-203-1-FPEPA,FAR-16-403-1-FPIF}.
  \item Used when scope is defined, cost estimates are reliable, and the Government wants to shift cost risk to the contractor~\autocite{edo-3-2-2-solicitation-prep-2025,FAR-16-103-Negotiating-Contract-Type,FAR-16-202-1-FFP}.
\end{itemize}
\begin{description}
  \item[\textbf{FFP.}] Price is not adjusted because of the contractor's cost experience; authorized changes or applicable clauses may still adjust price. The contractor assumes maximum risk and responsibility for costs and profit/loss~\autocite{edo-3-2-2-solicitation-prep-2025,FAR-16-202-1-FFP}.
  \item[\textbf{FPEPA.}] Price adjusted up or down based on specified contingencies (often economic indices)~\autocite{edo-3-2-2-solicitation-prep-2025,FAR-16-203-1-FPEPA}.
  \item[\textbf{FFP LOE.}] Fixed dollar amount for a specified level of effort over a stated period~\autocite{edo-3-2-2-solicitation-prep-2025,FAR-16-207-1-FFP-LOE}.
  \item[\textbf{FPIF.}] Final price equals final cost plus adjusted profit based on a shareline formula; profit varies with cost performance and is bounded by a ceiling price~\autocite{edo-3-2-2-solicitation-prep-2025,FAR-16-403-1-FPIF}.
  \item[\textbf{Fixed-price award fee.}] Fixed-price with award fee based on evaluation criteria when performance cannot be measured objectively~\autocite{edo-3-2-2-solicitation-prep-2025,FAR-16-404-FPAF}.
  \item[\textbf{Fixed-price with prospective price redetermination.}] Firm fixed price for an initial period with future price redetermination for later periods~\autocite{edo-3-2-2-solicitation-prep-2025}.
  \item[\textbf{Fixed-price with retroactive price redetermination.}] Fixed ceiling price with retroactive price redetermination after completion~\autocite{edo-3-2-2-solicitation-prep-2025}.
\end{description}
\begin{figure}[H]
  \centering
  \includegraphics[width=.78\linewidth]{FP_charts.png}
  \caption[Fixed-price contract types]{Fixed-price contract types and incentive geometry. \srcCite{edo-3-2-2-solicitation-prep-2025}.}
  \label{fig:fp_contract_types}
\end{figure}

\subsection{Fixed-Price Incentive (Firm Target) Example}
\begin{itemize}
  \item Given: target cost 100, target profit 10 (called target fee in the coursebook), target price 110, ceiling price 130, shareline 80/20~\autocite{edo-3-2-2-solicitation-prep-2025}.
  \item Signed cost variance = target cost minus actual cost (underrun positive; overrun negative)~\autocite{edo-3-2-2-solicitation-prep-2025}.
  \item Adjusted profit = target profit + (0.2 x signed cost variance)~\autocite{edo-3-2-2-solicitation-prep-2025}.
  \item Final price = the lesser of the ceiling price and actual cost plus adjusted profit~\autocite{edo-3-2-2-solicitation-prep-2025}.
\end{itemize}
\note{\ac{pta} = (ceiling price - target price) / Government share + target cost~\autocite{edo-3-2-2-solicitation-prep-2025}.}

\noindent New coursebook example: at actual cost 120, the signed variance is $100-120=-20$, adjusted profit is $10+0.2(-20)=6$, and final price is $120+6=126$, below the 130 ceiling~\autocite{edo-3-2-2-solicitation-prep-2025,FAR-16-403-1-FPIF}.

\subsection{Fixed-Price Incentive Practical Exercise}
\begin{itemize}
  \item Actual cost \$90,000; target cost \$100,000; target profit \$10,000; target price \$110,000.
  \item Shareline 80 Government / 20 contractor; ceiling price \$130,000.
  \item Underrun = \$10,000; adjusted profit = \$10,000 + (0.2 x \$10,000) = \$12,000.
  \item Final price = \$90,000 + \$12,000 = \$102,000.
\end{itemize}
~\autocite{edo-3-2-2-solicitation-prep-2025}

\subsection{Cost-Reimbursement Contracts}
\begin{itemize}
  \item Provide for payment of allowable incurred costs within a stated ceiling; the contractor exceeds the ceiling at its own risk unless the \ac{ko} authorizes an increase; applicable limitation-of-cost or limitation-of-funds clauses govern notice and continued performance~\autocite{edo-3-2-2-solicitation-prep-2025,FAR-16-301-1-Cost-Reimbursement}.
  \item Used when scope is uncertain, costs cannot be estimated accurately, an approved accounting system exists, and Government surveillance can ensure cost control~\autocite{edo-3-2-2-solicitation-prep-2025,FAR-16-103-Negotiating-Contract-Type,FAR-16-301-1-Cost-Reimbursement}.
  \item Place cost risk on the Government and can yield lower initial contract costs~\autocite{edo-3-2-2-solicitation-prep-2025,FAR-16-301-1-Cost-Reimbursement}.
\end{itemize}
\begin{description}
  \item[\textbf{CPFF.}] Allowable costs plus a fixed fee; fee can change only if scope changes~\autocite{edo-3-2-2-solicitation-prep-2025,FAR-16-306-CPFF}.
  \item[\textbf{CPIF.}] Allowable costs plus incentive fee tied to a formula based on total costs versus target costs; includes min and max fee~\autocite{edo-3-2-2-solicitation-prep-2025,FAR-16-405-1-CPIF}.
  \item[\textbf{CPAF.}] Allowable costs plus base fee (possibly zero) and award fee~\autocite{edo-3-2-2-solicitation-prep-2025,FAR-16-405-2-CPAF}.
  \item[\textbf{Cost.}] Allowable costs only; no fee~\autocite{edo-3-2-2-solicitation-prep-2025,FAR-16-302-Cost}.
  \item[\textbf{Cost sharing.}] Government pays an agreed portion of allowable costs; no fee~\autocite{edo-3-2-2-solicitation-prep-2025,FAR-16-303-Cost-Sharing}.
\end{description}
\begin{figure}[H]
  \centering
  \includegraphics[width=.78\linewidth]{CR_charts.png}
  \caption[Cost-reimbursement contract types]{Cost-reimbursement contract types and fee structures. \srcCite{edo-3-2-2-solicitation-prep-2025}.}
  \label{fig:cr_contract_types}
\end{figure}

\subsection{Cost-Plus-Incentive-Fee Example}
\begin{itemize}
  \item Given: target cost 1000, target fee 80, max fee 100, min fee 60, shareline 80/20~\autocite{edo-3-2-2-solicitation-prep-2025}.
  \item Signed cost variance = target cost minus actual cost (underrun positive; overrun negative)~\autocite{edo-3-2-2-solicitation-prep-2025}.
  \item Adjusted fee = target fee + (0.2 x signed cost variance), bounded by the minimum and maximum fee~\autocite{edo-3-2-2-solicitation-prep-2025}.
  \item Final price = actual cost + adjusted fee~\autocite{edo-3-2-2-solicitation-prep-2025}.
\end{itemize}

\noindent New coursebook example: at actual cost 900, the signed variance is $1000-900=100$, adjusted fee is $80+0.2(100)=100$ (the maximum fee), and final price is $900+100=1000$~\autocite{edo-3-2-2-solicitation-prep-2025,FAR-16-405-1-CPIF}.

\subsection{Incentives and Award Fees}
\begin{itemize}
  \item Incentives are driven by the requirement and can target cost, performance, schedule, award fee, options, or mixed \ac{clin} structures~\autocite{edo-3-2-2-solicitation-prep-2025,FAR-16-401-Incentive-Contracts}.
  \item Incentive contracts allow higher profit for meeting cost, performance, or schedule objectives when targets are clear and attainable~\autocite{edo-3-2-2-solicitation-prep-2025,FAR-16-401-Incentive-Contracts}.
  \item Most incentive contracts rely on cost incentives implemented via a shareline formula~\autocite{edo-3-2-2-solicitation-prep-2025,FAR-16-401-Incentive-Contracts}.
\end{itemize}
\begin{description}
  \item[\textbf{Special incentives.}] Tailored focus areas such as schedule milestones, key events, or drawing release counts; common in construction or design contracts~\autocite{edo-3-2-2-solicitation-prep-2025}.
  \item[\textbf{Award fee.}] Subjective criteria are used when performance cannot be measured objectively; award fee may be used in fixed-price or cost-plus contracts with base fee as applicable~\autocite{edo-3-2-2-solicitation-prep-2025}.
\end{description}
\note{Award fee criteria examples include systems engineering milestones, successful design reviews, test completion, and certifications~\autocite{edo-3-2-2-solicitation-prep-2025}.}

\subsection{Fixed-Price vs.\ Cost-Reimbursement Summary}
\begin{longtblr}[
    label   = {tab:fixed_vs_cost_reimb},
    caption = {Fixed-price vs.\ cost-reimbursement comparison},
    entry   = {Fixed-price vs.\ Cost-reimbursement},
    remark{\tSource} = {\tabCite{edo-3-2-2-solicitation-prep-2025,FAR-16-202-1-FFP,FAR-16-301-1-Cost-Reimbursement}}
  ]
  {colspec = {@{} ll @{}}}
  \toprule
  {Fixed-price} & {Cost-reimbursement} \\
  \midrule
  Promised outcome: acceptable goods or services & Promised outcome: best efforts \\
  Payment timing: after delivery (progress payments possible) & Payment timing: as costs are incurred \\
  Cost risk to contractor: high & Cost risk to contractor: low \\
  Cost risk to Government: low & Cost risk to Government: high \\
  \bottomrule
\end{longtblr}
\note{A contract may include multiple contract types through the use of \acp{clin}~\autocite{edo-3-2-2-solicitation-prep-2025,FAR-16-102-Contract-Type-Policies,FAR-16-103-Negotiating-Contract-Type}.}

\subsection{Contract Type by Acquisition Phase}
\begin{figure}[H]
  \centering
  \includegraphics[width=.95\linewidth]{contract_by_phase.png}
  \caption[Contract type by acquisition phase]{Contract type by phase of the acquisition process (notional coursebook model; legacy requirements labels retained). \srcCite{edo-3-2-2-solicitation-prep-2025}.}
  \label{fig:contract_by_phase}
\end{figure}
\note{Contract type selection is influenced by program risk and requirements definition~\autocite{edo-3-2-2-solicitation-prep-2025,FAR-16-103-Negotiating-Contract-Type}.}

\note{The figure's ICD/CDD-Validation/CDD labels describe the legacy classroom model. Apply current \ac{jfrp} and tailored requirements-document rules to the actual program; the figure does not prescribe current milestone documents.}

\subsection{Risk and Contract Types}
\begin{itemize}
  \item Cost risk shifts from Government to contractor as requirements become better defined and technical risk declines~\autocite{edo-3-2-2-solicitation-prep-2025,FAR-16-103-Negotiating-Contract-Type,FAR-16-202-1-FFP,FAR-16-301-1-Cost-Reimbursement}.
  \item Cost-reimbursement vehicles (\ac{cpff}, \ac{cpif}, \ac{cpaf}) align with higher technical uncertainty; incentive and fixed-price contracts align with defined requirements and determinable prices~\autocite{edo-3-2-2-solicitation-prep-2025,FAR-16-103-Negotiating-Contract-Type,FAR-16-301-1-Cost-Reimbursement,FAR-16-401-Incentive-Contracts}.
\end{itemize}

\subsection{Practical Exercise: Contract Type Selection}
\begin{itemize}
  \item Long-standing contractor relationship with a clear, well-defined scope and strong cost confidence.
  \item First-of-class depot-level overhaul for a new submarine class with uncertain cost drivers.
  \item Consider the contract type you would choose as the Government and as the contractor, tying the decision to risk.
\end{itemize}
~\srcCite{edo-3-2-2-solicitation-prep-2025}

\paragraph{Worked rationale.} Stable scope and reliable cost estimates support \ac{ffp}; the long-standing relationship alone does not justify that choice. A first-of-class overhaul with uncertain work may support cost reimbursement or an appropriate incentive arrangement if its prerequisites are met. The contractor prices the risk it accepts; the Government balances that price against uncertainty, oversight needs, and incentives~\autocite{edo-3-2-2-solicitation-prep-2025,FAR-16-103-Negotiating-Contract-Type}.

\subsection{Indefinite-Delivery Contracts}
\begin{itemize}
  \item Used when the Government has recurring needs but is uncertain about quantities or timing; reduces lead time for task or delivery orders~\autocite{edo-3-2-2-solicitation-prep-2025,FAR-16-501-2-Indefinite-Delivery,FAR-16-504-IDIQ}.
  \item Three types: definite-quantity (\ac{iddq}) commits to a specified quantity with deliveries scheduled by order; \ac{idiq} establishes a stated minimum and maximum with actual orders within those limits; a requirements contract covers the designated activity's actual purchase requirements during the stated period~\autocite{edo-3-2-2-solicitation-prep-2025,FAR-16-501-2-Indefinite-Delivery}.
\end{itemize}

\subsection{Other Contract Types}
\begin{description}
  \item[\textbf{T\&M.}] Used only when extent, duration, or costs cannot be estimated with reasonable confidence; pays labor at fixed hourly rates that include wages, overhead, \ac{ganda}, and profit, plus materials at actual cost, with a \ac{df} that no other contract type is suitable and a ceiling price the contractor exceeds at its own risk~\autocite{edo-3-2-2-solicitation-prep-2025,FAR-16-601-Time-Materials}.
  \item[\textbf{LH.}] Similar to \ac{tnm}, but without material charges~\autocite{edo-3-2-2-solicitation-prep-2025,FAR-16-602-Labor-Hour}.
  \item[\textbf{Letter contract (UCA).}] Preliminary instrument authorizing immediate performance with a price ceiling when definitive contract negotiations cannot be completed in time~\autocite{edo-3-2-2-solicitation-prep-2025,FAR-16-603-1-Letter-Contracts}.
\end{description}

\subsection{Socioeconomic Programs and Set-Asides}
\begin{itemize}
  \item Small business concerns are independently owned, not nationally dominant in their field, and meet \ac{sba} size standards. The board trap is treating ``small'' as one fixed employee or revenue threshold; current \ac{far}/\ac{sba} practice ties the applicable size standard to the solicitation's \ac{naics} code, and the size standard in effect on the solicitation date normally controls~\autocite{edo-3-2-2-solicitation-prep-2025,SBA-Size-Standards-2026,FAR-19-102-Size-Standards}.
  \item Socioeconomic programs include small business, woman-owned, veteran-owned, service-disabled veteran-owned, \ac{hubzone}, and \ac{sba} 8(a) certified firms~\autocite{edo-3-2-2-solicitation-prep-2025}.
  \item Set-asides reserve acquisitions exclusively for small businesses and, using the Oct~1,~2025 inflation-adjusted thresholds, actions above the \$15K micro-purchase threshold through the \$350K \ac{sat} are reserved for small business unless the \ac{ko} determines the rule-of-two/fair-market-price expectation is not met; above the \ac{sat}, set-asides are used when two or more responsible small businesses are expected and fair market prices are likely~\autocite{edo-3-2-2-solicitation-prep-2025,FAR-19-502-2-Set-Asides,AcquisitionGov-Threshold-Changes-2025}.
  \item Partial set-asides may be used for portions of multiple-award contracts, except construction~\autocite{edo-3-2-2-solicitation-prep-2025}.
  \item The \ac{sba} 8(a) program supports disadvantaged businesses with agency contracts and limited competition for up to nine years~\autocite{edo-3-2-2-solicitation-prep-2025}.
\end{itemize}

\subsection{Small Business Innovation Research}
\begin{itemize}
  \item \ac{sbir} enables small businesses to develop innovative technologies for Government use and later commercialization; \ac{sttr} is the parallel small-business/research-institution partnership path~\autocite{edo-3-2-2-solicitation-prep-2025,USC-638-SBIR,SBIR-Apply-2026}.
  \item Current statute requires participating agencies above the extramural research-and-development budget threshold to reserve not less than 3.2\% for \ac{sbir}; \ac{sttr} applies to agencies above the larger threshold and is generally 0.45\%. For \ac{dod}, the FY~2025--FY~2031 budget-calculation pilot uses at least 3.25\% of the two-year average \ac{rdte} extramural budget for \ac{sbir} and 0.46\% for \ac{sttr}~\autocite{USC-638-SBIR}.
  \item Three phases: Phase I proves technical merit and feasibility; Phase II continues the research and development; Phase III transitions/commercializes with non-\ac{sbir}/\ac{sttr} funds. Amounts and periods are solicitation-specific. Effective 16~October~2026, the SBA limits without an SBA waiver are \$335,937 for Phase~I and \$2,239,578 for Phase~II, including modifications; these are ceilings, not standard award amounts. Phase~III uses no \ac{sbir}/\ac{sttr} funding~\autocite{edo-3-2-2-solicitation-prep-2025,SBA-SBIR-Limits-20261016}.
\end{itemize}

\subsection{Other Transaction}
\begin{itemize}
  \item \ac{ot} instruments are not standard contracts, grants, or cooperative agreements and provide flexible business arrangements~\autocite{edo-3-2-2-solicitation-prep-2025}.
  \item Used for research and development, prototypes, or models to evaluate feasibility or utility; often funded with \ac{rdte}~\autocite{edo-3-2-2-solicitation-prep-2025}.
  \item Many Federal contract statutes do not apply, but procurement integrity rules and competitive practices still matter; \ac{ot} decisions can be protested at the \ac{gao} (limited) and the \ac{cofc}~\autocite{edo-3-2-2-solicitation-prep-2025}.
  \item Research-project \ac{ot} authority (10~U.S.C.\ \S~4021) and prototype-project / follow-on production authority (10~U.S.C.\ \S~4022) enable faster transitions when prototypes meet requirements~\autocite{USC-4021-OTA,USC-4022-OTF}.
  \item Evidence should include problem statements, solution briefs, prototype demo results, and transition plans to support production decisions~\autocite{DIU-Work-With-Us-2026}.
  \item \ac{gao} has noted that \ac{ot} use has expanded faster than data collection and outcome tracking; program teams should capture transition metrics early to support oversight~\autocite{GAO-20-84}.
\end{itemize}

\subsection{Data Rights Overview}
\begin{itemize}
  \item \ac{dod} policy emphasizes acquiring only the data rights essential to meet minimum Government needs~\autocite{edo-3-2-2-solicitation-prep-2025}.
  \item Categories include unlimited rights, government purpose rights, limited technical-data rights, restricted software rights, specifically negotiated licenses, special SBIR/STTR rights, and commercial-product/software licenses~\autocite{edo-3-2-2-solicitation-prep-2025}.
\end{itemize}
\begin{description}
  \item[\textbf{Unlimited rights.}] Government-funded development is one basis for unlimited rights in noncommercial technical data. Other categories, such as form-fit-function and operation, maintenance, installation, and training data (subject to clause exceptions), can also carry unlimited rights regardless of development funding. The Government may use or disclose such data for any purpose~\autocite{edo-3-2-2-solicitation-prep-2025,DFARS-252-227-7013-2026}.
  \item[\textbf{Limited rights.}] Generally applies to qualifying noncommercial technical data developed exclusively at private expense. Government use and disclosure outside Government are restricted, with clause-defined exceptions such as emergency repair/overhaul and covered support contractors under required safeguards~\autocite{edo-3-2-2-solicitation-prep-2025,DFARS-252-227-7013-2026}.
  \item[\textbf{Government purpose rights.}] Generally applies to mixed-funded development. Government purposes include competitive procurement during the rights period, with required restrictions on recipients; the rights normally become unlimited after five years unless a different period is negotiated~\autocite{edo-3-2-2-solicitation-prep-2025,DFARS-252-227-7013-2026}.
  \item[\textbf{Restricted rights.}] Applies to noncommercial computer software; Government use is limited (single machine use, minimum backup copies, transfer within Government, limited modification)~\autocite{edo-3-2-2-solicitation-prep-2025}.
\end{description}

\note{Development funding is a starting point, not the entire rights test. Distinguish technical data from computer software, commercial from noncommercial items, and the obligation to deliver data from the license to use it. Specifically negotiated and SBIR/STTR licenses require their own terms; inspect the applicable clause and markings before disclosure~\autocite{edo-3-2-2-solicitation-prep-2025,DFARS-252-227-7013-2026}.}

\subsection{Factors in Determining Data Rights}
\begin{itemize}
  \item Whether the item will be competitively acquired.
  \item Whether repair and overhaul will be contracted out or performed in-house.
  \item Whether parts are commercial items or require detailed drawings (build-to-spec or build-to-print).
\end{itemize}
\note{\ac{dod} policy requires procedures that are least intrusive on contractor economic interests as practicable~\autocite{edo-3-2-2-solicitation-prep-2025}.}

\subsection{Solicitation, Evaluation, and Award Phase}
\begin{itemize}
  \item Solicitations notify industry of a contracting need and its parameters~\autocite{edo-3-2-2-solicitation-prep-2025}.
  \item Method of contracting determines the solicitation type: \ac{ifb} for sealed bidding, \ac{rfp} for negotiated procurement~\autocite{edo-3-2-2-solicitation-prep-2025}.
  \item The \ac{ucf} is a common format for covered solicitations; procedure-specific exceptions apply. For negotiated acquisitions, current \ac{rfo} 15.102(b) states format exceptions and 15.109 describes the format. Submitted bids and proposals are offers~\autocite{edo-3-2-2-solicitation-prep-2025,FAR-Overhaul-Part15-2026}.
  \item Solicitations and pre-solicitation opportunities are publicized on \url{https://sam.gov}; proposed actions over \$25K are synopsized in the \ac{gpe}, and actions over \$20K through \$25K require public display or appropriate electronic notice unless an exception applies~\autocite{edo-3-2-2-solicitation-prep-2025,FAR-5-101-Disseminating-Information}.
\end{itemize}

\subsection{Uniform Contract Format Essentials}
\begin{itemize}
  \item Part I (The Schedule) includes Sections A--H; Part II contains contract clauses (Section I); Part III lists documents and attachments (Section J); Part IV covers representations and instructions (Sections K--M)~\autocite{edo-3-2-2-solicitation-prep-2025}.
  \item Section B provides supplies or services and prices/costs, with each item identified by a \ac{clin}~\autocite{edo-3-2-2-solicitation-prep-2025}.
  \item Section C provides descriptions, specifications, standards, and the \ac{sow}; it often references DD Form 1423 \ac{cdrl} requirements~\autocite{edo-3-2-2-solicitation-prep-2025}.
  \item Section J lists attachments, including \ac{cdrl} data requirements~\autocite{edo-3-2-2-solicitation-prep-2025}.
  \item Section L explains proposal submission requirements, organization, page limits, past performance data, and cost or pricing data~\autocite{edo-3-2-2-solicitation-prep-2025}.
  \item Section M states evaluation factors and their relative importance (e.g., price, past performance, technical, management, personnel, experience, schedule)~\autocite{edo-3-2-2-solicitation-prep-2025}.
\end{itemize}
\note{Parts~I--III form the resulting contract; Part~IV is retained in the contract file, with representations and certifications incorporated by reference as provided in the current format rules~\autocite{FAR-Overhaul-Part15-2026}.}
\note{Data deliverables add cost; balance reporting needs and \ac{cdrl} scope against program value~\autocite{edo-3-2-2-solicitation-prep-2025}.}

\subsection{Big-Picture Flow}
\begin{itemize}
  \item Requirements package \(\rightarrow\) contracting officer \(\rightarrow\) develop \ac{rfp} or \ac{ifb} \(\rightarrow\) advertise on \ac{sam} \(\rightarrow\) prospective offers~\autocite{edo-3-2-2-solicitation-prep-2025}.
\end{itemize}

\subsection{Quick Review}
\begin{itemize}
  \item Build a solicitation answer from need determination, requirements package, funds certification, market research, method of contracting, solicitation release, evaluation, award, administration, and closeout.
  \item Choose sealed bidding when the requirement is clear, price is the discriminator, discussions are unnecessary, and more than one responsible bid is expected.
  \item Choose negotiation when requirements, risk, price information, discussions, or best-value tradeoffs matter.
  \item Fixed-price fits defined requirements and shifts cost risk to the contractor; cost-reimbursement fits uncertain scope and shifts cost risk to the Government.
  \item \ac{idiq} vehicles reduce ordering lead time when recurring needs exist but exact quantities or timing are uncertain.
  \item In the uniform contract format, Section~L tells offerors how to propose and Section~M tells them how the Government will evaluate.
  \item Data rights, small-business strategy, \ac{sbir}/\ac{sttr}, \ac{ot}, and contract type should be decided early enough to shape the \ac{rfp}, not fixed after proposals arrive.
\end{itemize}

%====================
% End of file
%====================
\IfSubfilesClassLoaded{
  \RenewDocumentCommand{\entryname}{}{\textbf{\color{Modern} Acronym}}
  \RenewDocumentCommand{\descriptionname}{}{\textbf{\color{Modern} Definition}}
  \printnoidxglossary[
    type=\acronymtype,
    title=Acronyms,
    style=long-booktabs
  ]}{}
\end{document}
